\chapter{Machine-Learning Payload}

\section{What the Payload Does}

Lakitu carries a small camera and a dedicated AI chip pointed straight down
at the ground. Each capture cycle it takes a picture, runs it through a
small neural network on board, and reduces that picture to a single number:
how cloudy the sky looks, from 0\% (clear) to 100\% (overcast). Only that
number goes down the slow radio link -- the picture itself stays archived on
the SD card. The camera and its AI chip (a Coral Dev Board Micro with a
Google Edge TPU) form the ML payload; the STM32 OBC treats it as just
another sensor polled over a serial link.

The flight exists to prove exactly two things, and nothing more:

\begin{enumerate}
  \item \textbf{Shrink what has to be downlinked.} A raw frame is
  $\sim$50~KB, far more than the LoRa budget can carry; reducing it to a
  2-byte fraction on board is a $\sim$25{,}000:1 data reduction. Raw imagery
  is never downlinked, only archived.
  \item \textbf{Prove an Edge TPU survives and performs on a CubeSat-class
  platform.} A commercial, off-the-shelf AI accelerator never designed for
  spaceflight is asked to run continuous inference within a 1U CubeSat's
  power, thermal, mass, and bandwidth budget -- the claim is about
  integration, not about the cloud model being scientifically novel.
\end{enumerate}

\section{Proof-of-Concept Scope}

The PoC draws a hard boundary: the payload runs inference, produces
structured output, and hands it to the OBC over UART; packetization, LoRa
transmission, and ground reception belong to the OBC and TT\&C subsystems.
The engineering claim under test -- a commercial edge accelerator performing
useful on-board visual inference within a 1U CubeSat's envelope, delivering
structured output into the standard telemetry stream -- is treated as
falsifiable: Table~\ref{tab:payload-failure} lists the conditions under
which it fails.

\subsection{Minimum Demonstration}

The PoC counts as proven only if all four conditions below hold together.
The payload owns 1, 3, and 4 directly; 2 depends jointly on the OBC and
TT\&C link, so the OBC's SD archive (keyed by the same sequence number, SEQ,
as the downlink) is the fallback verification path if the downlink fails.

\begin{table}[H]
\centering
\begin{tabular}{lp{9cm}}
\toprule
\textbf{Criterion} & \textbf{What it means} \\
\midrule
Autonomous inference & Continuous output through the flight, no ground
intervention, at the configured cadence. \\
End-to-end downlink & Fraction records received and decoded correctly on the
ground via the OBC and TT\&C chain. \\
Within envelope & Power draw inside the EPS allocation; interior above
0\textdegree{}C; no thermal shutdown or throttling. \\
Verifiable against ground truth & Downlinked/archived cloud fraction agrees
with the archived source frame at a documented accuracy, post-flight. \\
\bottomrule
\end{tabular}
\caption{PoC minimum demonstration criteria}
\end{table}

\subsection{Failure Conditions}

\begin{table}[H]
\centering
\begin{tabular}{p{4.5cm}p{5cm}l}
\toprule
\textbf{Failure mode} & \textbf{Diagnosis} & \textbf{Attribution} \\
\midrule
No inference output produced & Coral, camera, or inference pipeline did not
run end-to-end & ML payload \\
Inference produced but not delivered to OBC & Coral output stage or UART
interface failed & ML payload \\
Inference delivered but not received on ground & OBC packetization, TT\&C
transmission, or ground reception failed & OBC / TT\&C / ground \\
Envelope exceeded (shutdown/throttle) & Power, thermal, or compute budget was
misjudged & ML payload \\
Output disagrees with ground truth & Incorrect inference, or data corrupted
in transit & ML payload / downstream \\
\bottomrule
\end{tabular}
\caption{PoC failure conditions and attribution}
\label{tab:payload-failure}
\end{table}

\subsection{Why a TPU at All}

The Edge TPU's value in this PoC is sustained inference throughput inside
the platform's power envelope, not enabling a task that would otherwise be
impossible on a CPU. Edge TPU inference latency for a MobileNet-class model
is on the order of milliseconds, so the model runs well inside each capture
cycle with margin to spare -- CPU-only inference at the same cadence would
draw more power per inference and erode the EPS budget instead.

\section{Implementation}

\subsection{Model: Two Candidates, One Deployed}

Two model formulations were trained on the same 95-Cloud / 38-Cloud
satellite imagery dataset: a \textbf{3-class classifier}
(\texttt{clear}/\texttt{overcast}/\texttt{partial}) and a
\textbf{continuous regressor} producing a single sigmoid cloud-fraction
output in $[0,1]$. The regressor was selected for flight -- strictly more
informative at the same output cost, and it keeps the telemetry field
(\texttt{fraction $\times$ 65535}) simple rather than needing a class enum.
The classifier pipeline remains in the repository but is not built into the
flight firmware.

Both share a MobileNet-class CNN backbone. The regressor is trained with
quantization-aware training (QAT) rather than plain post-training
quantization, then compiled for the Edge TPU; the deployed artifact is
committed directly under
\url{payload/coral/models/cloud_regressor_qat_int8_edgetpu.tflite} (an
exception to the repo's usual \texttt{*.tflite} ignore rule) so the firmware
builds without regenerating the model.

\subsection{Training Pipeline}

\begin{center}
\begin{tabular}{ll}
\toprule
\textbf{Stage} & \textbf{Artifact / tool} \\
\midrule
Raw imagery & 95-Cloud / 38-Cloud dataset (\texttt{payload/data/}, not
tracked) \\
Class dataset build & \texttt{training/build\_cloud\_classes.py} \\
Training & \texttt{training/notebooks/train\_cloud\_regressor.ipynb} (Keras)
\\
Quantization & QAT $\rightarrow$ int8 TFLite \\
Edge TPU compile & \texttt{edgetpu\_compiler} \\
Flight artifact & \texttt{coral/models/*\_edgetpu.tflite} (tracked) \\
\bottomrule
\end{tabular}
\end{center}

\subsection{Coral Firmware}

The flight firmware (\url{payload/coral/src/cloud_regressor.cc}) is an
out-of-tree \texttt{coralmicro} application (FreeRTOS on the M7 core). Per
cycle it captures a 224$\times$224 grayscale frame, invokes the int8 model
on the Edge TPU via an 8~MB tensor arena in SDRAM, and sends the result to
the OBC. Two bring-up findings are worth noting:

\begin{itemize}
  \item \textbf{UART6 is polled, not interrupt-driven.} The
  \texttt{coralmicro} fork moves the M7 debug console off LPUART6 onto
  LPUART7 to free the UART6 header for the OBC link; since the vendor's
  interrupt-driven driver still dispatches through per-instance state that
  is null for LPUART6 once the console moves, an LPUART6 interrupt would
  hard-fault, so the OBC link is polled I/O only.
  \item \textbf{SEQ survives reboots}, persisted to \texttt{/cloud\_seq}
  after every capture, so a mid-flight reset does not collide with
  already-archived frames on the OBC's SD card.
\end{itemize}

An unrecoverable initialisation failure triggers a self-reboot after a fixed
delay (\texttt{FatalRestart}) rather than hanging -- there is no one in orbit
to press reset. Debug-only tooling (HTTP/RPC image preview, burst capture)
is flagged in \texttt{docs/TODO.md} for removal before flight and does not
ship on the UART data path.

\subsection{Wire Interface}

Every cycle the payload sends one image packet to the OBC over UART6 at
115200~baud, 8N1: the full 224$\times$224 grayscale frame, the 2-byte cloud
fraction, and the sequence number, CRC-16 protected. The OBC archives the
frame to SD keyed by SEQ and forwards only the fraction record to the
downlink. The byte-level framing is the authoritative UART Protocol
Specification in \texttt{interface\_docs/}, not reproduced here.

\section{Thermal Envelope}

The payload carries its own requirement, distinct from the rest of the
platform: the Coral's interior must stay above 0\textdegree{}C throughout
the flight (MR-009 scopes this floor to the payload specifically). Passive
insulation alone cannot hold that floor through the night-cruise cold case;
a small thermostatic survival heater ($\sim$2--3~W) was evaluated as a
mitigation. Its cost is quantified in the EPS chapter: at 2.5~W and 30\%
duty the platform's power margin falls from $\sim$3.1$\times$ to
$\sim$1.9$\times$ -- still clear of the 1.5$\times$ floor.

Given the limited time before flight, the team chose passive-only thermal
management for this iteration and mitigated the risk by placement instead:
every high-internal-heat component (boost converter, batteries, Coral) sits
close together so the platform's own dissipation stays concentrated. This is
a managed risk -- the heater remains budgeted as a fallback, and the
freezer/vacuum session validates whether passive-only holds in practice.
